\emph{(i).} Let $e_{1}, \dots, e_{2m}$ be orthogonal generators of $\Cl(p,q)$ and $J = e_{1}\cdots e_{2m}$. Moving $e_{a}$ through $J$ passes $2m - 1$ generators that anticommute with it and one that commutes, so $e_{a}J = -Je_{a}$: $JvJ^{-1} = -v$ for every vector, and $\mathrm{Ad}(J)$ is the grade involution. Reversing the order of the $2m$ factors costs the sign $(-1)^{m(2m-1)} = (-1)^{m}$, so $J^{2} = (-1)^{m}\prod_{a}Q(e_{a}) = (-1)^{m+q}$. Changing the sign of one $Q(e_{a})$ changes $q$ by one. The listed cases follow; for the ten Majoranas $m = 5$, $q = 0$ and $J$ is the $\omega$ of \claimref{synthesis.order-four-elements-with-square-the-grading-and-the-majorana-element}(ii). The check computes the squares and the anticommutation for every signature with $2m \leq 8$ and for $(10,0)$.

\emph{(ii).} The identification is by an action of the group $\{1, \tau\}$ of the antipodal map exactly when the square of $(v, s \otimes w) \mapsto (-v, Js \otimes uw)$ is the identity, that is $(J \otimes u)^{2} = J^{2} \otimes u^{2} = 1$; the action on the sphere is free, so the quotient is a vector bundle. The Clifford bundle of a flat side is the quotient of $S^{n} \times \Cl(p,q)$ by $(v, a) \sim (-v, \hat{a})$, with $\hat{a} = JaJ^{-1}$ the grade involution, since $-1$ acts on the vectors of $\R^{2m}$ by $-1$. The fibrewise action $(v, a)\cdot(v, s \otimes w) = (v, as \otimes w)$ is compatible with the two identifications, because $\hat{a}\,Js \otimes uw = J(as) \otimes uw$. For $m + q$ even $J^{2} = 1$ and $u = 1$ serves; for $m + q$ odd $J^{2} = -1$ and the condition is $u^{2} = -1$.

\emph{(iii).} By the Whitney formula and $w(\shL) = 1 + x$ (Fact \claimref{synthesis.imported-whitney-product-formula-and-projective-space}), $w(E_{+})\,w(E_{-}) = (1 + x)^{2m}$. In degree one $w_{1}(E_{+}) + w_{1}(E_{-}) = 2m\,x = 0$; in degree two $w_{2}(E_{+}) + w_{2}(E_{-}) + w_{1}(E_{+})w_{1}(E_{-}) = \binom{2m}{2}x^{2} = m(2m - 1)x^{2} = m\,x^{2}$. On a flat side $t = p\,x = q\,x$ (as $p \equiv q$ modulo $2$), so the class is $(m + q)x^{2}$; directly, $\binom{p}{2} + \binom{q}{2} \equiv m + q \pmod 2$ for $p + q = 2m$, which the check verifies for $m \leq 6$. A definite side has $E_{-} = 0$, $t = 0$, class $m\,x^{2}$. On the far side of a wall with a definite physical side the dying line is the negative direction, $E_{-} = \shL$, so $t = x$ and the class is $(m + 1)x^{2}$. Exactly one of $m$, $m + 1$ is odd, and $x^{2} \neq 0$ in $H^{2}(\mathbb{RP}^{n}; \Z_{2})$ for $n \geq 2$. With (K) the side whose class is $x^{2}$ has no bundle of irreducible modules with structure group the spin group. For $m = 3$ the odd one is $m$ (the definite side) and for $m = 2$ it is $m + 1$ (the far side); the squares are those of (i).

\emph{(iv).} The check computes in $\Cl(6,2)$, with generators $e_{1}, \dots, e_{6}$ of square $+1$ and $f_{1}, f_{2}$ of square $-1$: $J_{6,0} = e_{1}\cdots e_{6}$ and $j = f_{1}f_{2}$ commute (each is even and they share no generator), $J_{6,0}^{2} = j^{2} = -1$, their product is the pseudoscalar, of square $+1$, and with $J_{5,1} = e_{1}\cdots e_{5}f_{1}$ and $E = e_{6}f_{2}$, $J_{5,1}E = -J_{6,0}\,j$, $E^{2} = 1$, $E$ anticommutes with $e_{6}$ and $f_{2}$ and commutes with the other six generators. Clause (4) of \claimref{synthesis.wall-rung-count-is-height-of-dying-line-class} is the statement that the far side's Clifford bundle, with the dying direction negative, tensored with the hyperbolic plane bundle $\Cl(\shL \otimes \R^{1,1})$, is the physical side's tensored with $\Cl(\shL \otimes \R^{0,2})$: here the physical side's dying direction is $e_{6}$, the far side's is $f_{1}$, the hyperbolic plane is spanned by $e_{6}$, $f_{2}$, and the jump plane by $f_{1}$, $f_{2}$. The monodromy $-1$ of $\shL$ acts on all eight directions by $-1$, and its lift is the pseudoscalar $J_{6,0}\,j = -J_{5,1}E$: on the far side $J_{5,1}$ times the plane's $E$, on the definite side the pair $(J_{6,0}, j)$. The unit bivector $j$ of the jump plane is fixed by $-1$ on the plane, $(-f_{1})(-f_{2}) = f_{1}f_{2}$, so it is a global section.

\emph{(v).} $w(\mathbb{RP}^{n}) = (1 + x)^{n+1}$ by the cited Fact, so $w_{1} = (n + 1)x$ and $w_{2} = \binom{n+1}{2}x^{2}$. $w_{1} = 0$ iff $n$ is odd; for $n = 2m - 1$, $\binom{2m}{2} \equiv m$, odd iff $n \equiv 1 \pmod 4$. For the order-four structure, $w(T\mathbb{RP}^{n} \oplus 2\shL) = (1 + x)^{n+3}$: $w_{1} = 0$ for $n$ odd and $w_{2} = w_{2}(\mathbb{RP}^{n}) + x^{2}$, which vanishes iff $w_{2}(\mathbb{RP}^{n}) = x^{2}$. With (L2) these are the stated existence statements; $x^{2} \neq 0$ needs $n \geq 2$. For the lens spaces the check computes the classes from the stored cell model instead of quoting them: in clause (3) of the cited theorem the diagonal of two odd cells has one term for each pair $0 \leq r < j \leq 3$, six terms, so $x^{2} = 6u = 0$ (one term, $x^{2} = u$, for the group of order two, which reproduces the classes of projective space above); and $w = (1 + u)^{k+1}$ in $\Z_{2}[u]/(u^{k+1})$. In words: clause (5) of \claimref{synthesis.cohomology-of-the-quarter-turn-lens-spaces-from-the-equivariant-cell-structure} gives $T\Sig_{k} \oplus \underline{\R} \cong (k + 1)\eta$ with $\eta$ an oriented plane bundle of Euler class $\pm y$, and clause (3) gives $u$ as the reduction of $y$, $H^{1}(\Sig_{k}; \Z_{2}) = \{0, x\}$ and $x^{2} = 0$. With the second half of (L2), $w(\eta) = 1 + u$, so $w_{1}(\Sig_{k}) = 0$ and $w_{2}(\Sig_{k}) = (k + 1)u$, which for $k = 2$ is $u \neq 0$. An order-four structure with a line of class $a$ needs $w_{2} = a^{2} = 0$. So $L^{5}(4)$ has neither.

\emph{(vi).} The torus is generated by $Y$, $T_{3}$ and two colour generators, all diagonal on the occupation states, with $6Y$, $2T_{3}$ and the colour weights integers. An element is $\exp(2\pi I(\tfrac{a}{12}\,6Y + \tfrac{b}{4}\,2T_{3} + \dots))$; those with square $1$ on the sixteen have all phases half-integers of a turn, and it suffices to search $a$ modulo $12$ and the others modulo $4$. The check enumerates them: sixteen distinct sign patterns on the sixteen states, the trivial one and fifteen that reverse exactly eight states; none reverses the empty state, the neutral singlet, on which all four weights vanish. An element $h$ of the connected group with $h^{2} = 1$ on the sixteen lies in a maximal torus, all of which are conjugate (classical, not derived here), and conjugation by a group element permutes eigenvalues with their multiplicities, so $h$ reverses none or eight states and fixes the singlet, which is the trivial representation. $\omega$ acts on the sixteen by one unit $-I$ (\claimref{synthesis.order-four-elements-with-square-the-grading-and-the-majorana-element}(ii)) and commutes with the torus, so $(\omega g)^{2} = -g^{2} = -1$, and the signed count of $\omega g$ is $-16 + 2f$ on sixteen components and $-15 + 2f$ on fifteen, $f \in \{0, 8\}$ the number of reversed states: $-16$, $0$ and $-15$, $1$. The neutral singlet is the state whose presence a gauge algebra inside $\uu(5)$ leaves free (\claimref{synthesis.neutral-singlet-bit-decided-by-gauge-algebra-leaving-u5}). The counts on the class of $\omega$ are $-16$ and $0$ (respectively $-15$ and $1$), both attained, so the differences have greatest common divisor $16$ and no larger modulus makes the count a function of the class. The classes modulo $16$ and the statement on generations are arithmetic. $X = 5 - 2N$ takes the values $5$, $1$, $-3$ on the $1$, $10$, $5$ states with $N = 0, 2, 4$: $\sum X = 5 + 10 - 15 = 0$, $\sum X^{3} = 125 + 10 - 135 = 0$; without the empty state $-5$ and $-125$; with the full state ($X = -5$) in its place the linear sum is $-10$. These are the mixed gravitational and cubic sums of Fact \claimref{standard-model.imported-anomaly-coefficients} for the charge $X$; the check recomputes them.

\emph{The Reading} uses (i) to (vi), the count of \claimref{synthesis.spectral-asymmetry-on-real-projective-five-space-from-monogenic-polynomials}(iv), (v) as (H1), and the hypothesis (H2) it names; nothing in it is a proved statement beyond these.
